\chapter{Special Functions from Equations and Integrals}
\label{ch:special-equations}

A differential equation is useful only after its chosen solution has been
specified and its computation bounded.  A recurrence supplies local data;
the fundamental matrix moves that data along certified paths.  Bessel
and hypergeometric functions illustrate the two parts, and a Bessel zero
shows how a special-function value can itself become an inverse computation.

\section{Bessel functions of integer order}
For an integer $m\ge0$, define
\[
 J_m(z)=\sum_{k\ge0}\frac{(-1)^k(z/2)^{2k+m}}{k!(k+m)!}.
\]
Every coefficient is rational.  On $|z|\le R$, the absolute terms have
successive ratio
\[
 \frac{R^2}{4(k+1)(k+m+1)}.
\]
This decreases to zero; after it is at most $1/2$, the tail is at most
twice the first omitted majorant term.  Multiplying by any fixed falling
factorial gives a derivative-tail sequence with the same eventual property.
Consequently these are entire function and derivative computations.  To
obtain a geometric coefficient chart of any chosen radius, evaluate the
absolute series majorant at a strictly larger radius and bound each
coefficient by that weighted sum.

Coefficient comparison proves
\[
 z^2J_m''+zJ_m'+(z^2-m^2)J_m=0,
 \qquad J_0'=-J_1.
\]
For example the coefficient of degree $2k+m$ in the equation is
$4k(k+m)a_k+a_{k-1}=0$ for $k\ge1$; the lowest coefficient is free
and is set to $1/(2^m m!)$.  This both verifies the equation and records
which local solution was selected.  Standard notation and the equation
are in \href{https://dlmf.nist.gov/10.2}{DLMF, Section 10.2}.

On a path avoiding zero, the state $(J_m,J_m')$ satisfies
\[
 Y'=\begin{pmatrix}0&1\\m^2/z^2-1&-1/z\end{pmatrix}Y.
\]
Start at a nonzero point with values from the entire series.  Pull back
each rational polygonal path segment $z=p+t(q-p)$ to $0\le t\le1$:
its coefficient matrix is $(q-p)A(p+t(q-p))$.  A certified positive
separation from zero bounds all entries and supplies their continuity.
The Peano--Baker construction then propagates the state.  Integral-equation
uniqueness proves agreement with the series.  It is unnecessary, and
incorrect, to apply the bounded-coefficient evolution theorem at $z=0$;
the series handled that point separately.

\section{A second Bessel computation}
The circle functions give the bounded integral
\[
 B(z)=\int_0^1 E(izC(t))\,dt.
\]
The factorial majorant is uniform for $0\le t\le1$ on every bounded
$z$-disk.  Thus it can be integrated term by term.  Symmetry cancels odd
powers, and finite trigonometric integration gives
\[
 \int_0^1 C(t)^{2k}\,dt=\frac{\binom{2k}{k}}{4^k}.
\]
For completeness, differentiating $S(t)C(t)^{2k-1}$ and integrating its
zero endpoint difference yields the recurrence
$M_k=(2k-1)M_{k-1}/(2k)$, with $M_0=1$.
The product has the uniform increment estimates already proved for sine
and cosine, so this use of the FTC is justified.
Substitution into the exponential coefficients shows $B=J_0$.
For real $z$ its imaginary part cancels by $t\mapsto1-t$; hence
\[
 J_0(z)=\int_0^1\cos(zC(t))\,dt,
\]
where $\cos w=(E(iw)+E(-iw))/2$.  This is the normalized Bessel integral
of \href{https://dlmf.nist.gov/10.9}{DLMF, Section 10.9}, proved here
by finite moments and a uniform tail.

\section{Computing a zero without a general root theorem}
\label{sec:bessel-zero}
Exact alternating-series bounds give
\[
 J_0(2)>\frac29>0,\qquad J_0(3)<-\frac{4199}{16384}<0.
\]
The first bound uses terms through $k=3$; the second uses terms through
$k=4$.  Although the first absolute term at $3$ need not decrease, the
remaining tail after these cutoffs does decrease and has the indicated
sign.  At $0\le x\le3$, writing $q=x^2$, the tail in $J_1$ after its
four displayed terms gives
\[
 J_1(x)\ge\frac x2
       \left(1-\frac q8+\frac{q^2}{192}-\frac{q^3}{9216}\right).
\]
The polynomial in parentheses decreases on $[0,9]$: its derivative is
$-1/8+q/96-q^2/3072<0$.  Its value at $9$ is $223/1024$.
Therefore $J_1(x)\ge223/1024$ on $[2,3]$, and the certified FTC gives
\[
 J_0(u)-J_0(v)\ge\frac{223}{1024}(v-u)
 \quad(2\le u<v\le3).
\]
The endpoint signs and this separation make the trisection inverse
construction of Chapter~\ref{ch:foundations} terminate.  They compute a
unique zero in $[2,3]$ as nested rational intervals.  Also
$J_0(x)\ge1-x^2/4>0$ for $0\le x<2$, so this is its first positive
zero.  The root is a new number computation, not a symbolic placeholder
chosen by a completeness theorem.

\section{A recurrence with a finite convergence disk}
For rational parameters $a,b,c$ with $c\notin\{0,-1,-2,\ldots\}$, put
\[
 h_0=1,\qquad
 (n+1)(n+c)h_{n+1}=(n+a)(n+b)h_n,
 \qquad F(z)=\sum_{n\ge0}h_nz^n.
\]
If an upper parameter makes a numerator zero, the recurrence terminates
and gives a polynomial.  Otherwise, for a chosen rational $r<1$, bound
the term ratio by
\[
 r\frac{(n+|a|)(n+|b|)}{(n+1)(n-|c|)}\qquad(n>|c|).
\]
Choose rational $r<\theta<1$.  After clearing the positive denominator,
the assertion that this ratio is at most $\theta$ is a quadratic
polynomial inequality in $n$ with positive leading coefficient.  Choose
an integer $N$ large enough that its expansion in $n=N+j$ has nonnegative
coefficients.  That finite check proves the ratio bound for every $j\ge0$.
The tail after $N$ is then bounded by its first omitted majorant term
divided by $1-\theta$.  Derivative tails are handled by the same method.

To supply a chart on $|z|<R<1$, apply this ratio argument at $R$ and
bound the finitely many earlier weighted coefficients.  This gives a
rational $M$ with $|h_n|\le M R^{-n}$ for all $n$.
Coefficient comparison now proves
\[
 \boxed{z(1-z)F''+[c-(a+b+1)z]F'-abF=0.}
\]
The recurrence and $h_0=1$ identify the local solution, denoted
${}_2F_1(a,b;c;z)$; see
\href{https://dlmf.nist.gov/15.2}{DLMF, Section 15.2} and
\href{https://dlmf.nist.gov/15.10}{Section 15.10} for the standard notation.

\section{Continuation uses the equation, not repeated recentering alone}
Away from $0$ and $1$ the preceding equation is a first-order system with
\[
 A(z)=\begin{pmatrix}
 0&1\\
 ab/[z(1-z)]&((a+b+1)z-c)/[z(1-z)]
 \end{pmatrix}.
\]
The convergent series gives starting data at a nonzero point inside the
unit disk.  A finite path with certified separation from $0,1$ gives
bounded coefficient computations on each pulled-back segment; the
fundamental matrix propagates the state along the path.

It also supplies fresh analytic charts.  Here is a finite majorant
argument for any rational-coefficient system at an ordinary point.
Suppose locally $A(z_0+w)=\sum A_jw^j$ with
$\|A_j\|\le M R^{-j}$, obtained by polynomial and reciprocal charts.
The solution coefficients satisfy
\[
 (n+1)y_{n+1}=\sum_{j=0}^n A_jy_{n-j}.
\]
Choose $0<r\le R/2$ with $2Mr\le1$, and a rational
$C\ge\|y_0\|$.  Induction gives
\[
 \|y_n\|\le Cr^{-n}:
 \quad \|y_{n+1}\|\le\frac{2MC}{n+1}r^{-n}
                                      \le Cr^{-n-1}.
\]
On $|w|\le r/2$ the local series and all requested derivative tails
are therefore effective.  Its integral equation identifies it with the
propagated state.  Coefficient separation from the singularities allows
this bound to be renewed at successive centers.

Distinct paths may produce distinct branches.  The finite overlap and
path-deformation certificates of Chapter~\ref{ch:complex-paths} say when
they agree; there is no assumption of global single-valuedness.

\section{The beta integral identifies the same local function}
For positive rational $b,c-b$ and rational $a$, define on $|z|<1$
\[
 H(z)=\frac1{B(b,c-b)}\int_0^1
       t^{b-1}(1-t)^{c-b-1}(1-zt)^{-a}\,dt.
\]
Use the local logarithm to define the last factor.  Its coefficient
recurrence is $(n+1)d_{n+1}=(n+a)d_n$, $d_0=1$, obtained from
$(1-w)D(1-w)^{-a}=a(1-w)^{-a}$.  The ratio estimate just used supplies
uniform tails on smaller disks.  The positive beta weight provides the
endpoint majorant, so the parameter-integral theorem applies.

The beta--gamma identity and gamma recurrence give
\[
 \frac{B(b+n,c-b)}{B(b,c-b)}
       =\frac{b(b+1)\cdots(b+n-1)}{c(c+1)\cdots(c+n-1)}.
\]
Thus integrating the binomial coefficients yields exactly the $h_n$
above, and $H=F$.  This proves equivalence between an improper integral,
a convergent series, and a propagated ODE computation.
In particular the elliptic integral developed later has
\[
 K(m)=\frac\pi2\,{}_2F_1(1/2,1/2;1;m).
\]
The later geometric substitution gives its independent bounded quadrature;
the coefficient identity here explains why its equation has the displayed
hypergeometric form.
